% Pattern Info

\tikzset{
pattern size/.store in=\mcSize,
pattern size = 5pt,
pattern thickness/.store in=\mcThickness,
pattern thickness = 0.3pt,
pattern radius/.store in=\mcRadius,
pattern radius = 1pt}
\makeatletter
\pgfutil@ifundefined{pgf@pattern@name@_bgrd8j6uq}{
\pgfdeclarepatternformonly[\mcThickness,\mcSize]{_bgrd8j6uq}
{\pgfqpoint{0pt}{0pt}}
{\pgfpoint{\mcSize}{\mcSize}}
{\pgfpoint{\mcSize}{\mcSize}}
{
\pgfsetcolor{\tikz@pattern@color}
\pgfsetlinewidth{\mcThickness}
\pgfpathmoveto{\pgfqpoint{0pt}{\mcSize}}
\pgfpathlineto{\pgfpoint{\mcSize+\mcThickness}{-\mcThickness}}
\pgfpathmoveto{\pgfqpoint{0pt}{0pt}}
\pgfpathlineto{\pgfpoint{\mcSize+\mcThickness}{\mcSize+\mcThickness}}
\pgfusepath{stroke}
}}
\makeatother
\tikzset{every picture/.style={line width=0.75pt}} %set default line width to 0.75pt

\begin{tikzpicture}[x=0.75pt,y=0.75pt,yscale=-1,xscale=1]
%uncomment if require: \path (0,724); %set diagram left start at 0, and has height of 724

%Shape: Ellipse [id:dp12751974195854354]
\draw  [pattern=_bgrd8j6uq,pattern size=12pt,pattern thickness=0.22pt,pattern radius=0pt, pattern color={rgb, 255:red, 215; green, 215; blue, 215}] (100,371.8) .. controls (100,290.72) and (196.75,225) .. (316.1,225) .. controls (435.45,225) and (532.2,290.72) .. (532.2,371.8) .. controls (532.2,452.88) and (435.45,518.6) .. (316.1,518.6) .. controls (196.75,518.6) and (100,452.88) .. (100,371.8) -- cycle ;
%Shape: Parabola [id:dp7097132516162088]
\draw   (584.7,459.6) .. controls (406.53,205.73) and (228.37,205.73) .. (50.2,459.6) ;
%Shape: Parabola [id:dp18585049318765634]
\draw   (48.85,270) .. controls (227.02,523.87) and (405.18,523.87) .. (583.35,270) ;
%Shape: Ellipse [id:dp6935180515795655]
\draw  [fill={rgb, 255:red, 255; green, 255; blue, 255 }  ,fill opacity=1 ] (236.5,371.8) .. controls (236.5,346.78) and (272.14,326.5) .. (316.1,326.5) .. controls (360.06,326.5) and (395.7,346.78) .. (395.7,371.8) .. controls (395.7,396.82) and (360.06,417.1) .. (316.1,417.1) .. controls (272.14,417.1) and (236.5,396.82) .. (236.5,371.8) -- cycle ;


% 原图 6.5 中沿开口朝上的抛物线设置的法向箭头。
\draw[->] (156,392)--(169,375);
\draw[->] (193,420)--(205,399);
\draw[->] (231,442)--(240,418);
\draw[->] (268,455)--(273,429);
\draw[->] (300,460)--(301,432);
\draw[->] (332,460)--(331,432);
\draw[->] (364,454)--(359,428);
\draw[->] (401,442)--(391,418);
\draw[->] (438,420)--(426,400);
\draw[->] (477,392)--(464,375);

% Text Node
%Straight Lines [id:da6462481686425137]
\draw    (317.45,269.2) ;
\draw [shift={(317.45,269.2)}, rotate = 180] [fill={rgb, 255:red, 0; green, 0; blue, 0 }  ][line width=0.08]  [draw opacity=0] (8.93,-4.29) -- (0,0) -- (8.93,4.29) -- cycle    ;
%Straight Lines [id:da8521475132911274]
\draw    (316.1,460.4) ;
\draw [shift={(316.1,460.4)}, rotate = 180] [fill={rgb, 255:red, 0; green, 0; blue, 0 }  ][line width=0.08]  [draw opacity=0] (8.93,-4.29) -- (0,0) -- (8.93,4.29) -- cycle    ;

\draw (146,357.4) node [anchor=north west][inner sep=0.75pt]    {$a$};
% Text Node
\draw (481,357.4) node [anchor=north west][inner sep=0.75pt]    {$b$};
% Text Node
\draw (311,466.4) node [anchor=north west][inner sep=0.75pt]    {$C_{0}$};
% Text Node
\draw (300,244.4) node [anchor=north west][inner sep=0.75pt]    {$C'_{0}$};
% Text Node
\draw (194,266.4) node [anchor=north west][inner sep=0.75pt]    {$N$};
% Text Node
\draw (325,337.4) node [anchor=north west][inner sep=0.75pt]    {$S$};

\end{tikzpicture}
